% Section 10.3, from lectures/L10/appendix/c_subroutines_and_stack.md.

\section{Subrutiner och stacken}
\label{c10:app:c}

\subsection{Samma kod på flera ställen}
\label{c10:c1}
Programmen har vuxit. Blinkprogrammet i \secref{c10:a10} har en fördröjning på nio rader, och ett
trafikljus behöver fyra fördröjningar av olika längd. Att skriva av samma nio rader fyra gånger
fungerar, men det är fyra ställen att rätta när något är fel, och fyra uppsättningar etiketter som
måste heta olika.

Hex-till-ASCII-programmet i \secref{c10:a3} har samma problem i liten skala: samma fem rader två
gånger.

Lösningen är en \textbf{subrutin}: en bit kod som skrivs en gång, får ett namn, och
\textbf{anropas} från så många ställen som behövs. När subrutinen är klar fortsätter programmet där
anropet gjordes. Andra programspråk kallar samma sak en funktion eller en metod.

\subsection{\code{rcall} och \code{ret}}
\label{c10:c2}
Två instruktioner räcker:

\begin{booktable}{@{}p{7.4em}Lc@{}}
  \thead{Instruktion} & \thead{Gör} & \thead{Cykler} \\
  \midrule
  \code{rcall delay_ms} & Sparar adressen till instruktionen efter \code{rcall}, och hoppar till
    \code{delay_ms}. & 3 \\
  \code{ret} & Hoppar tillbaka till den sparade adressen. & 4 \\
\end{booktable}

Adressen som sparas kallas \textbf{returadressen}: den visar vart subrutinen ska återvända. Men var
sparas den? Inte i något av de 32 registren, utan på \textbf{stacken}, ett område i toppen av
dataminnet.

\needspace{10\baselineskip}
\begin{asmcode}
loop:
    ...
    rcall delay_ms                  ; Save the address of rjmp loop, and jump.
    rjmp loop                       ; ret comes back here.

delay_ms:
    ...                             ; The subroutine's work.
    ret                             ; Jump back to the saved address.
\end{asmcode}

\noindent Ett anrop kostar alltså 3 + 4 = 7 cykler innan subrutinen har gjort något alls. Det är
lite, men det är skälet till att en fördröjning i en subrutin blir några cykler längre än samma
loop skriven direkt i programmet.

\subsection{Stacken och stackpekaren}
\label{c10:c3}
\textbf{Stacken} är en del av SRAM som används som en trave: det som läggs dit sist tas bort först.
Den börjar i toppen av dataminnet, på adress \code{RAMEND} = \code{0x08FF}, och växer
\textbf{nedåt}, mot lägre adresser.

\textbf{Stackpekaren}, \code{SP}, är ett 16-bitars register som pekar på den första lediga byten i
stacken. Den består av två I/O-register, \code{SPH} (hög byte) och \code{SPL} (låg byte).

Så här går ett anrop till, byte för byte:
\begin{enumerate}
  \item Före \code{rcall} pekar \code{SP} på \code{0x08FF}, och stacken är tom.
  \item \code{rcall} skriver returadressens låga byte på \code{0x08FF} och den höga på
    \code{0x08FE}, och flyttar \code{SP} två steg nedåt, till \code{0x08FD}. Sedan hoppar den.
  \item \code{ret} flyttar \code{SP} två steg uppåt, läser returadressen, och hoppar dit.
\end{enumerate}

\bookfigure{avr_call_stack}{Stacken före \code{rcall}, i subrutinen och efter
\code{ret}.}{c10:fig:call_stack}

Lägg märke till det sista: \code{ret} \textbf{raderar ingenting}. Bytena ligger kvar i minnet, men
\code{SP} pekar förbi dem, och nästa anrop skriver över dem.

\subsubsection{Att initiera stackpekaren}
Ett program som använder subrutiner börjar med att peka \code{SP} på toppen av SRAM:

\needspace{7\baselineskip}
\begin{asmcode}
main:
    ldi r16, high(RAMEND)           ; Point the stack pointer at the top of SRAM,
    out SPH, r16                    ; high byte first ...
    ldi r16, low(RAMEND)
    out SPL, r16                    ; ... then the low byte. SP = 0x08FF.
\end{asmcode}

\noindent På ATmega328P gör hårdvaran faktiskt redan detta vid varje reset, så i våra program
ändrar de fyra raderna ingenting. De skrivs ändå, av två skäl: många andra AVR-kretsar startar med
\code{SP} = 0, och ett program som hoppar in i ditt program utan reset, till exempel en bootloader,
kan ha lämnat \code{SP} var som helst. Fyra rader är billig försäkring, och de gör programmet
begripligt för den som läser det.

\subsection{En subrutin med parameter: \code{delay_ms}}
\label{c10:c4}
En fördröjning som alltid är lika lång är till begränsad nytta. Bättre är en subrutin som får veta
hur länge den ska vänta. Värdet den får kallas en \textbf{parameter}, och i kursen lämnas den i
ett register före anropet, här \code{r24}:

\needspace{4\baselineskip}
\begin{asmcode}
    ldi r24, 250                    ; Wait 250 ms.
    rcall delay_ms
\end{asmcode}

\noindent Hela subrutinen, från \code{blink_subroutine.asm}:

\begin{asmcode}
;--------------------------------------------------------------------------------
; delay_ms: Waits r24 milliseconds (1-255) at 16 MHz.
;
; Each turn of the outer loop takes exactly 16 000 cycles, 1 ms. With the rcall, the
; pushes, the pops and the ret, a call takes 16 000 x r24 + 18 cycles.
;--------------------------------------------------------------------------------
delay_ms:
    push r24                        ; Save the registers this subroutine changes.
    push r26
    push r27
delay_ms_outer:
    ldi r26, low(3999)              ; X = r27:r26 = 3999.
    ldi r27, high(3999)
delay_ms_inner:
    sbiw r26, 1                     ; 4 cycles per turn, 3 the last time.
    brne delay_ms_inner
    dec r24                         ; One more millisecond done.
    brne delay_ms_outer
    pop r27                         ; Restore them, in the opposite order.
    pop r26
    pop r24
    ret                             ; Back to the instruction after the rcall.
\end{asmcode}

\noindent Uppbyggnaden är samma som sekundfördröjningen i \secref{c10:a10}: en 16-bitars inre loop
i en yttre loop. Den yttre loopen räknar ned \code{r24}, en millisekund per varv. Ett varv tar
\mbox{\code{4 · 3999 + 4 = 16 000}} cykler, precis 1~ms vid 16~MHz.

Resten av cyklerna är subrutinens egna:

\begin{narrowtable}{@{}lc@{}}
  \thead{Del} & \thead{Cykler} \\
  \midrule
  \code{rcall} & 3 \\
  tre \code{push} & 6 \\
  \code{r24} varv om 16\,000 cykler, minus 1 för sista \code{brne} & 16\,000 · \code{r24} -- 1 \\
  tre \code{pop} & 6 \\
  \code{ret} & 4 \\
\end{narrowtable}

\noindent Summan är \textbf{16\,000 · \code{r24} + 18} cykler. Med \code{r24} = 250 blir det
4\,000\,018 cykler, och det är också vad simulatorn mäter mellan \code{rcall} och instruktionen
efter. De 18 extra cyklerna tar drygt en mikrosekund, en miljondels sekund. Ingen människa ser
skillnad, men det är bra att veta var de kommer ifrån.

Varför börjar och slutar subrutinen med \code{push} och \code{pop}? Det är ämnet för de två nästa
avsnitten.

\subsection{\code{push} och \code{pop}}
\label{c10:c5}
Stacken kan användas till mer än returadresser. Två instruktioner lägger dit och hämtar godtyckliga
register:

\begin{booktable}{@{}lLc@{}}
  \thead{Instruktion} & \thead{Gör} & \thead{Cykler} \\
  \midrule
  \code{push r16} & Skriver \code{r16} där \code{SP} pekar, och flyttar \code{SP} ett steg nedåt. &
    2 \\
  \code{pop r16} & Flyttar \code{SP} ett steg uppåt, och läser byten där till \code{r16}. & 2 \\
\end{booktable}

Stacken är \textbf{sist in, först ut}. Programmet \code{push_pop.asm} visar vad det betyder:

\needspace{8\baselineskip}
\begin{asmcode}
    ldi r16, 0x11
    ldi r17, 0x22
    push r16                        ; 0x08FF = 0x11, SP = 0x08FE.
    push r17                        ; 0x08FE = 0x22, SP = 0x08FD.
    pop r16                         ; r16 = 0x22 (the last one pushed), SP = 0x08FE.
    pop r17                         ; r17 = 0x11, SP = 0x08FF.
\end{asmcode}

\bookfigure{l10_push_pop}{Stacken när \code{r16} och \code{r17} läggs dit och hämtas tillbaka i
samma ordning.}{c10:fig:push_pop}

Registren har bytt värden, eftersom \code{pop r16} får det som lades dit sist. Vill man ha tillbaka
samma värden i samma register ska de hämtas i \textbf{omvänd ordning}: \code{push r16},
\code{push r17}, sedan \code{pop r17}, \code{pop r16}.

\subsection{Kursens regel: spara det du ändrar}
\label{c10:c6}
En subrutin använder register för sitt arbete. \code{delay_ms} räknar med \code{r24}, \code{r26}
och \code{r27}. Men programmet som anropar den kan ha något viktigt i just de registren, och det
vet inte subrutinen.

Kursens regel är därför enkel:

\begin{aside}
\textbf{En subrutin sparar varje register den ändrar med \code{push} när den börjar, och återställer
dem med \code{pop} i omvänd ordning innan \code{ret}.} Undantaget är ett register som subrutinen är
till för att lämna ett svar i; det ska stå i kommentaren ovanför subrutinen.
\end{aside}

Med regeln behöver den som anropar en subrutin aldrig undra om ett register har ändrats. Efter
\code{rcall delay_ms} innehåller \code{r24} fortfarande 250, och därför kan samma parameter
användas igen utan att laddas om.

Så här ser stacken ut mitt i \code{delay_ms}, efter anropet och de tre \code{push}-instruktionerna,
i \code{blink_subroutine.asm}. \code{rcall} ligger på adress \code{0x000B} i programminnet, så
returadressen är \code{0x000C}:

\bookfigure{l10_stack_in_delay}{Stacken i \code{blink_subroutine.asm}: före \code{rcall}, efter
\code{rcall}, efter tre \code{push} och efter \code{pop} och \code{ret}.}{c10:fig:stack_in_delay}

Fem byte är upptagna: två för returadressen och tre för de sparade registren. \code{SP} pekar på
\code{0x08FA}. Tre \code{pop} tar bort registren i omvänd ordning, \code{ret} tar returadressen,
och \code{SP} är tillbaka på \code{0x08FF}.

\subsubsection{Vad som händer om man slarvar}
Anta att \code{delay_ms} gör tre \code{push} men bara två \code{pop}. Då ligger en byte kvar ovanpå
returadressen när \code{ret} körs, och \code{ret} läser den byten som en del av adressen.
Programmet hoppar till en adress som ingen har bestämt, någonstans i programminnet, och fortsätter
därifrån. Ingenting varnar. Ofta hamnar programmet i tomt flashminne, kör igenom det, och börjar om
från adress 0, så att det ser ut som om kortet startar om av sig självt.

Det är därför regeln säger \textbf{varje} \code{push} har sin \code{pop}, i omvänd ordning. Ett fel
här syns inte när man läser koden rad för rad, men det syns direkt i simulatorn: stega fram till
\code{ret} och titta på \code{SP}. Står den inte där den stod när subrutinen började, stämmer inte
stacken.

\subsection{Subrutiner som anropar subrutiner}
\label{c10:c7}
En subrutin kan anropa en annan. Varje anrop lägger en ny returadress på stacken, och varje
\code{ret} tar bort den senaste, så programmet hittar alltid tillbaka i rätt ordning.

\needspace{9\baselineskip}
\begin{asmcode}
delay_500ms:
    push r24                        ; Save the register this subroutine changes.
    ldi r24, 250
    rcall delay_ms                  ; 250 ms. delay_ms saves r24, so it is still 250 ...
    rcall delay_ms                  ; ... and this is another 250 ms.
    pop r24
    ret
\end{asmcode}

\noindent När huvudprogrammet anropar \code{delay_500ms}, som anropar \code{delay_ms}, ligger det
på stacken:
\begin{itemize}
  \item returadressen till huvudprogrammet, 2 byte;
  \item \code{r24}, sparat av \code{delay_500ms}, 1 byte;
  \item returadressen till \code{delay_500ms}, 2 byte;
  \item \code{r24}, \code{r26} och \code{r27}, sparade av \code{delay_ms}, 3 byte.
\end{itemize}

Åtta byte, och \code{SP} står på \code{0x08FF - 8 = 0x08F7}. Det är den djupaste punkten i
programmet.

Lägg också märke till att det andra anropet av \code{delay_ms} fungerar utan att \code{r24} laddas
om. Det är regeln i \secref{c10:c6} som gör det möjligt.

\subsection{Hur mycket stack finns det?}
\label{c10:c8}
SRAM är 2048 byte. Programmen i kursen använder som mest ett tiotal byte stack, så det finns gott
om plats. Men det finns ingen spärr: om stacken skulle växa ned i området där programmet har egna
data, till exempel texten i \secref{c10:a5}, skriver den över dem utan varning.

Därför är det värt att räkna efter, på samma sätt som i \secref{c10:c7}: två byte per anrop som
pågår samtidigt, och en byte per sparat register. Summan är stackens största djup.

\subsection{Hex-till-ASCII som subrutin}
\label{c10:c9}
Nu kan programmet i \secref{c10:a3} skrivas utan de dubbla raderna. Omvandlingen av en nibble blir
en subrutin som tar värdet 0--15 i \code{r24} och lämnar tecknet i samma register. Det är just det
undantag regeln i \secref{c10:c6} tillåter: \code{r24} är till för svaret, och kommentaren säger
det.

\needspace{12\baselineskip}
\begin{asmcode}
main:
    ldi r16, high(RAMEND)           ; SP = RAMEND.
    out SPH, r16
    ldi r16, low(RAMEND)
    out SPL, r16

    ldi r16, 0x3C                   ; The byte to convert.
    mov r24, r16
    swap r24
    andi r24, 0x0F                  ; The high nibble, 0x03.
    rcall nibble_to_ascii
    mov r20, r24                    ; r20 = '3'.
    mov r24, r16
    andi r24, 0x0F                  ; The low nibble, 0x0C.
    rcall nibble_to_ascii
    mov r21, r24                    ; r21 = 'C'.
end:
    rjmp end

;--------------------------------------------------------------------------------
; nibble_to_ascii: Turns the value 0-15 in r24 into its hexadecimal character.
; Returns the character in r24; changes no other register.
;--------------------------------------------------------------------------------
nibble_to_ascii:
    cpi r24, 10                     ; A digit 0-9, or a letter A-F?
    brlo nibble_digit
    subi r24, -7                    ; A-F: 7 more, so that 10 lands on 'A'.
nibble_digit:
    subi r24, -0x30                 ; Add the code for '0'.
    ret
\end{asmcode}

\noindent Subrutinen ändrar flaggorna, med \code{cpi} och \code{subi}. Det gör nästan varje
subrutin, och det är därför huvudprogrammet aldrig ska räkna med att flaggorna är orörda efter ett
anrop. Ett villkorligt hopp ska alltid komma direkt efter den jämförelse det bygger på.
